\scanpage{32}
\mspage{31}
\mspage{55 ter}

\uline{Démonstration.} Le cas a) n'est autre que 5.10 et a été mis pour mémoire.

b) \deleted{\unsure{Le cas $\ell\ne p$ n'est autre que le critère galoisien de ré}}\inserted{La première assertion est une conséquence \deleted{\unsure{immé}} triviale de} a) et des définitions. Pour prouver la dernière assertion dans b), on peut \deleted{donc} supposer que $\ell\ne p$, mais alors grâce à a) il suffit de prendre pour K' une extension finie séparable de K correspondant à un sous-groupe ouvert $\pi'$ de $\pi=\operatorname{Gal}(\bar K/K)$ dont l'intersection avec le groupe d'inertie I est contenue dans le \deleted{\unsure{stabi}} noyau de la représentation de I sur $T_\ell(A_K(\bar K))$, lequel est \deleted{\unsure{fi}} ouvert par hypothèse (cf. 5.12.1).

c) On a déjà signalé que $1)\Rightarrow2)$ (5.12.2), il est trivial que $2)\Rightarrow3)$, reste à prouver que $3)\Rightarrow1)$. Pour $\ell\ne p$, ce n'est autre que \inserted{le} critère galoisien de réduction semi-stable 3.5. Pour $\ell=p$, on peut procéder ainsi. Choisissons un nombre premier $\ell'\ne p$, et appliquons à celui-ci les considérations de 4.1. \deleted{On est donc ramené à prouver que la représentation} Supposons S strictement local, ce qui est loisible (3.4.0), prenons une extension galoisienne finie K' de K, de groupe G, telle que $A_{K'}$ ait réduction semi-stable. Il faut prouver que la représentation de G sur les trois quotients intervenant dans (4.1.1) est triviale, ou ce qui revient au même, que la \deleted{somme des traces} représentation somme est triviale, ou encore, que la trace de cette dernière est ``triviale'', i.e. constante . Or nous avons vu dans 4.3 a) que cette fonction trace est indépendante \unsure{da|de} choix de $\ell'\ne p$, et dans 6.\unsure{4} plus bas nous indiquons une méthode de calcul de cette fonction en \deleted{fonction}\inserted{termes} de $T_p(A_K)$ également. Ce calcul donne trivialement une fonction constante dans le cas où $T_p(A_K)$ a bonne réduction virtuelle, et cela achève la démonstration.

\uline{Remarque} 5.13.1. Il semble plausible que si $X_\eta$ est un \inserted{pro-$\ell$-} groupe de Barsotti-Tate sur $\eta$ qui a bonne réduction virtuelle sur S, et bonne réduction virtuelle essentielle d'échelon n, alors il a bonne réduction virtuelle d'échelon n. (C'est \deleted{\unsure{ce qui avait été utilisé, pour}} vrai pour $\ell\ne p$.)
